\section{Related work}\label{sec:related}

We discuss prior work along three questions: which sparse nonconvex problems
are tractable, how grid and vertex bounds are computed and certified, and when
the work of a method is independent of the accuracy. The last paragraphs
cover exact output and the extensions.

\paragraph{Tractability of sparse nonconvex quadratic programs.}
For binary quadratic and pseudo-Boolean optimization, bounded treewidth gives
polynomial algorithms by nonserial dynamic programming
\cite{BerteleBrioschi1972,CramaHansenJaumard1990,Dechter1999} and exact
lift-and-project or linear formulations of size exponential only in the width
\cite{BienstockOzbay2004,WainwrightJordan2004,Laurent2009}; sparse moment and
sum-of-squares hierarchies exploit the same structure for polynomial problems
\cite{WakiEtAl2006,Lasserre2006}. For continuous box QP the situation differs.
Del Pia and Khajavirad give an exact strongly polynomial dynamic program on
forests and prove strong NP-hardness at treewidth two, even with bounded
integral coefficients, and strong NP-hardness of quartic minimization on paths
\cite[Theorems~1--3]{DelPiaKhajavirad2026}. Their tractable classes of
unbounded treewidth, and those of Khajavirad \cite{Khajavirad2026PolyBox},
take variables with nonpositive diagonal to be binary;
Remark~\ref{rem:cv-instances} compares Khajavirad's elimination of the
continuous components with our value factors.
Exact dynamic programs with piecewise-quadratic messages, for box QP on
forests \cite{DelPiaKhajavirad2026} and for total-variation problems on trees
\cite{KolmogorovPockRolinek2016,KuricAhmetspahicPock2024}, are compared with
our grids in Section~\ref{sec:limits-messages}. Box QP is also
polynomial when the rank of the Hessian is fixed
\cite{HladikCernyRada2021}. For convex quadratics with indicator variables,
Bhathena et al.\ give an $O(n^2)$ parametric dynamic program on trees
\cite{BhathenaEtAl2026Trees} and, on graphs of bounded treewidth, one whose
running time is linear in $n$ for fixed width, margin, volume growth and
condition numbers \cite{BhathenaEtAl2026Graphs}. This is the closest
treewidth-and-conditioning result we know; its objective is convex, with
indicator constraints (see also \cite{BienstockChen2024}). For integer QP,
bounded treewidth gives fixed-parameter tractability together with bounded
domains \cite{EibenEtAl2019}; parameterized by the number of variables,
bounded coefficients are also needed \cite{Lokshtanov2015}, and without them
the problem is W[1]-hard \cite{Herrmann2026}; see \cite{GanianOrdyniak2018}
for integer linear programming. Approximation schemes polynomial in
$1/\varepsilon$ exist in fixed dimension
\cite{DeLoeraEtAl2008,HildebrandWeismantelZemmer2016}, for sparse polynomial
problems with constraints \cite[Theorem~4]{BienstockMunoz2018}, for QP with a
fixed number of negative eigenvalues \cite{Vavasis1992}, for mixed-integer
QP (MIQP) with fixed rank and a fixed number of integer variables
\cite{DelPia2023}, and for MIQP with a fixed number of negative eigenvalues
and of integer variables
\cite{DelPia2026Jacobi}. Several of these works explain why the dependence on
$1/\varepsilon$ cannot be made polylogarithmic in their settings unless
P${}={}$NP \cite{BienstockMunoz2018,Vavasis1992,DelPia2026Jacobi}, from the
NP-hardness of Subset Sum or of quadratic programming with one negative
eigenvalue \cite{PardalosVavasis1991}. Running times that depend on a
condition number of the instance are standard in numerical computation
\cite{BurgisserCucker2013}. To our knowledge, no earlier result gives a
running time of the form $f(p,\kappa)\poly(I+q)$, or an exact algorithm in
time $f(p,\kappa)\poly(I)$, for nonconvex or mixed-integer box QPs.

\paragraph{Grid and vertex bounds, filtering and certificates.}
Lower bounds from function values on a grid are classical for Lipschitz
functions \cite{Piyavskii1972,Shubert1972,Nesterov2004}; with second-order
information the error is quadratic in the mesh
\cite{BreimanCutler1993,DeKlerkLasserreLaurentSun2017}. Separable quadratic
perturbations that convexify a function ($\alpha$BB) have maximum separation
$\frac14\sum_i\alpha_i(u_i-\ell_i)^2$
\cite{MaranasFloudas1994,AndroulakisMaranasFloudas1995,AdjimanEtAl1998};
perturbations of the opposite sign make a function edge-concave, so that its
envelope is vertex polyhedral
\cite{Tardella2004,MeyerFloudas2005,Hasan2018,MisenerFloudas2012}. Bajaj and
Hasan bound a function on a box from below by its values at the $2^n$
vertices, using an upper bound on the diagonal Hessian entries
\cite{BajajHasan2020}. Piecewise-linear and sawtooth relaxations of squares
apply the same second-order error termwise inside a mixed-integer program (MIP)
\cite{BeachHildebrandHuchette2022,BeachEtAl2024,DongLuo2018}.
Proposition~\ref{prop:cellwise} observes that, with unary corrections at the
grid nodes, the best such vertex bound over all cells of a product grid is
the minimum of a factorized function, computed by one tree dynamic program
instead of an enumeration of cells.

Removing a region whose relaxation value exceeds an incumbent is
optimality-based bound tightening, probing and branch-and-reduce
\cite{RyooSahinidis1996,BelottiEtAl2009,GleixnerEtAl2017,PuranikSahinidis2017};
cost-based domain filtering \cite{FocacciLodiMilano1999} and dead-end
elimination \cite{DesmetEtAl1992,GainzaRobertsDonald2012} apply the same rule
in constraint programming and protein design. Exact min-marginals of
tree-structured problems are computed by two passes of message passing
\cite{Dawid1992,AjiMcEliece2000,WainwrightJaakkolaWillsky2005,KohliTorr2006},
and thresholding them prunes states in structured prediction cascades
\cite{WeissTaskar2010}. Our filter is an instance of this principle; its role
is that, under growth, the retained coordinate hulls shrink with the mesh.
Exact rational mixed-integer programming \cite{CookKochSteffyWolter2013} and
independently checkable branch-and-bound certificates
\cite{CheungGleixnerSteffy2017} guard MIP results against floating-point
error; our path certificate follows the logic of the latter, with a tree that
is a path. Interval methods and safe bounds control rounding errors in
global optimization and MIP \cite{Hansen1980,HansenWalster2004,
NeumaierShcherbina2004}; our certificates use exact rational arithmetic
instead.

\paragraph{Accuracy-independent work under growth.}
In branch-and-bound with second-order convergent bounds, the number of boxes
that remain near a nondegenerate minimizer is independent of the tolerance
(the cluster problem); this count is local and asymptotic, and it can be
exponential in the dimension and in the conditioning of the Hessian
\cite{DuKearfott1994,Neumaier2004,WechsungSchaberBarton2014,KannanBarton2017}.
Lemmas~\ref{lem:inv} and~\ref{lem:states} are a global, non-asymptotic form of
this phenomenon for product grids on a tree decomposition. Because filtering
is coordinatewise, the exponential dependence falls on the bag size rather
than on the dimension, and the condition $8\bar\kappa\theta^2\le1$ plays the
role of the bound on the second-order prefactor. Optimistic optimization
attains exponential rates when the growth order at the optimum matches the
assumed smoothness \cite{Munos2011}, and for Lipschitz functions whose
near-optimal sets are small, branch-and-bound needs $O(\log(1/\varepsilon))$
function values \cite{Perevozchikov1990}. Proximity and scaling give
algorithms logarithmic in the domain size for linear and separable convex
integer programs \cite{CookGerardsSchrijverTardos1986,HochbaumShanthikumar1990,
GranotSkorinKapov1990,HunkenschroderEtAl2026,EisenbrandEtAl2025}. Quadratic
growth and error bounds give linear convergence of local methods
\cite{LuoTseng1993,KarimiNutiniSchmidt2016,DrusvyatskiyLewis2018,
NecoaraNesterovGlineur2019}, and restarts remove the need to know the growth
constant \cite{RouletDAspremont2020}; our capped trials play the same role.

\paragraph{Adaptive discretization.}
Coarse-to-fine dynamic programming replaces groups of states by superstates
with optimistic costs, so that each coarse problem is a lower bound, refines
along the current optimal path, and converges to a unique optimum
\cite{Raphael2001,FelzenszwalbMcAllester2007}. Interval convex max-product
bounds potentials on interval cells to obtain discrete lower bounds for
continuous Markov random fields \cite{PengEtAl2011}. Interval dynamic
programming for optimal power flow on trees returns approximately feasible
superoptimal solutions with a number of bound evaluations polynomial in
$1/\varepsilon$ \cite[Theorem~2]{DvijothamEtAl2017}. Continuous distributed
constraint optimization \cite{HoangEtAl2020,TroullinosEtAl2022},
variable-resolution discretization in control \cite{MunosMoore2002}, corridor
methods \cite{HeidariEtAl1971} and adaptive two-level dynamic programming
\cite{LeaverFayKuhlmanSnoeyink2005} also refine discretizations adaptively.
In mixed-integer nonlinear programming, adaptive multivariate partitioning
refines piecewise relaxations around the current relaxation solution and
combines them with optimality-based bound tightening, with no bound on the
number of partition points needed for a given accuracy
\cite{NagarajanEtAl2019}, and a primal heuristic re-centers discretizations
of fixed size on the previous solution \cite{GupteKosterKuhnke2022}. None of
these methods certifies a given accuracy with a number of states independent
of that accuracy.

\paragraph{Exact rational output.}
A rational QP that attains its minimum has a global minimizer of polynomial
encoding length \cite{Vavasis1990,DelPiaDeyMolinaro2017}; the
stationary-polytope argument of Section~\ref{sec:exact} follows Vavasis, and
Remark~\ref{rem:np} discusses the relation to NP membership. Rounding an
approximation to an exact rational is classical
\cite[Section~5.1]{GrotschelLovaszSchrijver1988}, \cite{KozlovTarasovKhachiyan1980},
and uniqueness of a QP minimizer implies quadratic growth by second-order
optimality theory for quadratic programs \cite{Contesse1980}. Growth towards
the whole optimal set of a quadratic program over a polytope follows from the
error bound of Luo and Sturm \cite[Theorem~3.3]{LuoSturm2000}. For polynomials
of degree three or more, irrational optimizers and encoding-size phenomena
limit exact output \cite{BienstockDelPiaHildebrand2023}.

\paragraph{Value functions and multiparametric programming.}
Eliminating private variables through their optimal response is a standard
form of decomposition, as in Benders decomposition
\cite{Benders1962,Geoffrion1972} and condensing in model predictive control
\cite{BockPlitt1984,JerezKerriganConstantinides2012,Axehill2015}. When the
eliminated block has active bounds, multiparametric quadratic programming
gives piecewise-affine optimizers and piecewise-quadratic values on polyhedral
critical regions \cite{BemporadEtAl2002,TondelJohansenBemporad2003}, also for
positive semidefinite Hessians with nonunique optimizers
\cite{SpjotvoldTondelJohansen2007,PatrinosSarimveis2011}. Value functions of
parameters that enter the objective linearly over a fixed feasible set are
concave \cite{MangasarianRosen1964,FiaccoKyparisis1986}, and infima preserve
semiconcavity \cite{CannarsaSinestrari2004}. We use these facts for a different
purpose: an upper coordinate curvature bound for the reduced objective, valid
globally, so that the corrected grid remains a certified bound. Low negative
inertia with convex recourse is treated by spatial branching in
\cite{Vavasis1992,LuoBaiLimPeng2019}.

\paragraph{Cuts and submodularity.}
Binary quadratic functions with nonpositive pair coefficients are minimized by
one minimum cut \cite{Ivanescu1965,PicardRatliff1975,BorosHammer2002,
KolmogorovZabih2004}, after sign switching whenever the signed interaction
graph is balanced \cite{Harary1953,HammerHansenSimeone1984,Padberg1989}.
Variables with nonpositive diagonal can be fixed to endpoints
\cite{Rosenberg1972,DelPiaKhajavirad2026,Khajavirad2026PolyBox}, so continuous
quadratics that are submodular and concave in each coordinate reduce to such
cuts \cite{BurerNatarajanWillemsen2026v3}. Pairwise problems submodular with
respect to label orders reduce to cuts on threshold encodings
\cite{SchlesingerFlach2006,Ishikawa2003,Hochbaum2001}. For general continuous
submodular quadratics, Burer, Natarajan and Willemsen show that a
semidefinite relaxation is exact up to dimension three, give a
four-variable gap and leave the complexity open from dimension four on
\cite{BurerNatarajanWillemsen2026v3}; a sign-restricted case is exact
\cite{KimKojima2003}, and discretization methods
\cite{Bach2019,Bach2018Isotonic,AxelrodLiuSidford2020} are polynomial in the
inverse accuracy. Under point growth, Theorem~\ref{thm:balanced} is instead
polynomial in the input length, $\kappa$ and $\log(1/\varepsilon)$ for
balanced box quadratics; since $\kappa$ can be exponential in the input
length, it does not settle that question. Partial minimization preserves submodularity
\cite{Topkis1978,Topkis1998}; this underlies the mixed continuous--discrete
methods of \cite{GomezHan2025,BuntonTabuada2022}, with certificates by convex
combinations of greedy extreme bases
\cite{Edmonds1970,Cunningham1985,IwataFleischerFujishige2001}.

\paragraph{Constraints, optimal sets and limits.}
Integrality of totally unimodular fibers is the Hoffman--Kruskal theorem
\cite{HoffmanKruskal1956,Schrijver1986}; marginal-preserving rounding under hard
constraints is classical in approximation algorithms
\cite{GandhiEtAl2006,ChekuriVondrakZenklusen2010}. Remark~\ref{rem:tu-bm}
compares our extension to totally unimodular coupling constraints with the
linear programming approximations of Bienstock and Mu\~noz for constrained
polynomial problems \cite{BienstockMunoz2018}. Section~\ref{sec:endpointset}
relates our description of the optimal set to Rosenberg's face criterion and
to tree reparameterizations \cite{Rosenberg1972,WainwrightJaakkolaWillsky2005,Werner2007}. Diagonal Lagrangian
certificates for box quadratics are given in
\cite{BeckTeboulle2000,JeyakumarRubinovWu2006,LiWuQuan2015}, and exactness of
relaxations of box QP is characterized in \cite{QiuYildirim2024}. Sparse
second-order cone and semidefinite formulations are exact under graph
conditions \cite{SojoudiLavaei2014,DeyKhajavirad2026,Khajavirad2026SparseSDP};
chordal sparsity guarantees positive semidefinite completions
\cite{GroneEtAl1984,VandenbergheAndersen2015}. Marginals that agree on
separators glue into a joint measure on tree-structured families
\cite{Vorobev1962,WainwrightJordan2008}, but locally consistent moments, as in
lift-and-project relaxations, need not come from a joint measure
\cite{SheraliAdams1990,WainwrightJordan2008}; complete parametric descriptions
can have exponentially many pieces \cite{Murty1980,MairalYu2012,GartnerJaggiMaria2012}.
